\documentclass[11pt,a4paper]{article}
\usepackage[utf8]{inputenc}
\usepackage[T1]{fontenc}
\usepackage[brazilian]{babel}
\usepackage{geometry}
\usepackage{graphicx}
\usepackage{booktabs}
\usepackage{float}
\usepackage{enumitem}
\usepackage{amsmath}
\usepackage{xcolor}
\usepackage{tcolorbox}
\geometry{margin=2.5cm}
\setlength{\parindent}{0pt}
\setlength{\parskip}{0.5em}
\tcbuselibrary{skins,breakable}
\newtcolorbox{nota}[1][]{colback=yellow!10,colframe=orange!60!black,fonttitle=\bfseries,title=#1,breakable}

\begin{document}
\begin{center}
  
\includegraphics[width=0.16\textwidth]{../../../figures_global/unb.png}\\[0.3cm]
  {\large\textbf{GABARITO DO PROFESSOR -- USO RESTRITO}}\\[0.2cm]
  {\large FGA0092 -- Princípios de Comunicação para Engenharia}\\[0.2cm]
  \rule{12cm}{0.8pt}\\[0.2cm]
  {\Large\textbf{Prática 07: Ruído em Demoduladores}}\\[0.2cm]
  \rule{12cm}{0.8pt}
\end{center}

\section*{Gráfico esperado}
\begin{center}
  \includegraphics[width=0.88\textwidth]{grafico_esperado.png}
\end{center}
\textit{O gráfico acima foi gerado pelo script \texttt{gerar\_grafico.py} desta pasta.}

\section*{Respostas esperadas}

\begin{enumerate}[label=\textbf{Q\arabic*.}, leftmargin=*, itemsep=1em]

  \item \textbf{Como medir o SNR de um sinal no GNU Radio usando blocos disponíveis?}

    \textbf{Resposta esperada:} O SNR pode ser estimado usando o bloco \texttt{RMS},
    que mede a potência RMS do sinal. Com dois blocos RMS -- um no sinal antes do
    canal (sinal limpo) e outro no ruído isolado -- calcula-se
    $\text{SNR} = P_\text{sinal} / P_\text{ruído}$ (razão de potências). Outra abordagem
    visual utiliza a amplitude máxima do sinal versus o piso de ruído no osciloscópio
    (\texttt{QT GUI Time Sink}).

    \begin{nota}[Erros comuns]
      Medir amplitude pico em vez de potência RMS, resultando em SNR incorreto por
      um fator de $\sqrt{2}$ para sinais senoidais (diferença de $3$\,dB).
    \end{nota}

  \item \textbf{Como a qualidade do sinal demodulado AM DSB-SC se degrada com o aumento do ruído?}

    \textbf{Resposta esperada:} O demodulador AM DSB-SC é linear: a degradação do SNR
    de saída é proporcional à do canal ($\text{SNR}_o = \text{SNR}_i$ em escala linear).
    A relação é aproximadamente linear: dobrar o ruído reduz o SNR pela metade.
    A degradação visual começa quando \texttt{noise\_sigma} $\approx$ amplitude do sinal
    ($\text{SNR} \approx 0$\,dB). Para \texttt{noise\_sigma}$= 0{,}05$ com amplitude 1,0,
    $\text{SNR} \approx 26$\,dB; a qualidade ainda é alta.

    \begin{nota}[Erros comuns]
      Afirmar que o AM tem ``degradação exponencial'' com o ruído, ou não identificar
      o limiar de degradação visual como próximo de SNR$\,= 0$\,dB.
    \end{nota}

  \item \textbf{Como o FM se degrada com ruído, e o que são os ``cliques'' observados na saída?}

    \textbf{Resposta esperada:} O FM tem comportamento não-linear: acima do limiar
    ($\text{SNR}_i > \sim 10$\,dB), a qualidade de saída é muito superior ao AM.
    Abaixo do limiar, o discriminador de FM colapsa abruptamente, gerando ``cliques''
    de ruído impulsivo de alta amplitude na saída. Este fenômeno é claramente visível
    no GNU Radio como picos abruptos na saída do \texttt{WBFM Receive}, e é característica
    distinta do FM em comparação com a degradação gradual do AM.

    \begin{nota}[Erros comuns]
      Não mencionar o fenômeno de limiar do FM, ou descrevê-lo como uma degradação
      gradual semelhante ao AM, ignorando o caráter impulsivo dos cliques.
    \end{nota}

  \item \textbf{Para $\beta=5$ (\texttt{max\_dev}$= 10000$\,Hz, $f_m = 2000$\,Hz), qual o ganho teórico do FM sobre o AM? Como isso se reflete experimentalmente?}

    \textbf{Resposta esperada:} Ganho teórico:
    $3\beta^2(\beta+1) = 3 \times 25 \times 6 = 450$. Isso significa que o FM suporta
    \texttt{noise\_sigma} aproximadamente $\sqrt{450} \approx 21$ vezes maior que o AM
    para a mesma qualidade de saída. Experimentalmente, o FM deve mostrar vantagem para
    \texttt{noise\_sigma} entre 0,05 e 0,5. O custo é $\approx 11\times$ mais largura
    de banda ($B_\text{FM} = 22$\,kHz vs.\ $B_\text{AM} = 4$\,kHz).

    \begin{nota}[Erros comuns]
      Comparar AM e FM sem normalizar pela mesma potência transmitida ou mesma largura
      de banda, o que favorece artificialmente um dos dois sistemas.
    \end{nota}

  \item \textbf{O que é uma comparação justa entre AM e FM em termos de ruído? Quais cuidados são necessários?}

    \textbf{Resposta esperada:} Para comparação justa: (1) mesma potência do sinal de
    mensagem; (2) mesmo SNR de entrada no canal; (3) como o FM usa largura de banda
    maior, a potência de ruído total no receptor FM é maior -- para SNR idêntico em
    densidade espectral ($N_0$), o ruído total FM é $(B_{FM}/B_{AM})$ vezes maior.
    Comparar pela mesma \texttt{noise\_sigma} absoluta favorece artificialmente o FM,
    pois não contabiliza a diferença de largura de banda.

    \begin{nota}[Erros comuns]
      Ignorar que o FM recebe mais potência de ruído total devido à sua maior largura
      de banda, tornando a comparação injusta quando se usa a mesma \texttt{noise\_sigma}.
    \end{nota}

\end{enumerate}

\section*{Critérios de correção}
\begin{tabular}{ll}
  \toprule
  Critério & Pontuação \\
  \midrule
  Q1 -- Uso correto de potência RMS; distinção de amplitude pico vs.\ RMS & 1,0 ponto \\
  Q2 -- Degradação linear do AM DSB-SC com SNR e identificação do limiar   & 1,0 ponto \\
  Q3 -- Fenômeno de limiar do FM e caráter impulsivo dos cliques            & 1,0 ponto \\
  Q4 -- Cálculo do ganho teórico e conversão para fator de noise\_sigma     & 1,0 ponto \\
  Q5 -- Critérios de comparação justa AM vs.\ FM com correção de banda      & 1,0 ponto \\
  \midrule
  \textbf{Total} & \textbf{5,0 pontos} \\
  \bottomrule
\end{tabular}

\bigskip
\begin{nota}[Nota de correção geral]
  Priorizar coerência física e evidência experimental (capturas de tela, tabelas).
  Respostas sem evidência visual devem ter desconto. Não penalizar pequenos erros
  de cálculo se o raciocínio físico estiver correto.
\end{nota}

\end{document}
